\documentclass[12pt,a4paper]{article}
\usepackage[utf8]{inputenc}
\usepackage[T2A]{fontenc}
\usepackage{lmodern}
\usepackage{setspace}
\onehalfspacing
\usepackage[ukrainian]{babel}
\usepackage{geometry}
\geometry{margin=2.5cm}

\usepackage{sectsty}
\usepackage{titlesec}
\usepackage{enumitem}
\usepackage{hyperref}

% Налаштування заголовків
\titleformat{\section}{\large\bfseries}{\thesection.}{1em}{}
\titleformat{\subsection}{\normalsize\bfseries}{\thesubsection.}{1em}{}

\hypersetup{
    colorlinks=true,
    linkcolor=blue,
    citecolor=blue,
    urlcolor=blue
}

\title{Інтелект як інструмент самообману: когнітивні та соціальні механізми}
\author{Y. M. Volynsky}
\date{}

\begin{document}

\maketitle
\vspace{1em}
\begin{center}
\small
Матеріали до міждисциплінарного семінару\\
«Психологія ідеологій та когнітивні спотворення»\\
травень 2024
\end{center}
\begin{abstract}
У статті аналізується, як інтелектуальні здібності індивіда можуть використовуватися не для досягнення епістемічної точності, а для стабілізації самооцінки, зниження когнітивного дисонансу та захисту вже сформованих позицій. Розглядаються взаємодія мотивованого мислення, раціоналізації, інтелектуалізації та селективної рефлексії. Показано, що вищий рівень когнітивної складності часто не зменшує, а ускладнює і зміцнює форми самообману.
\end{abstract}

\section{Вступ}
Традиційна раціоналістична модель мислення припускає, що зростання освіти та когнітивних здібностей веде до зменшення помилок і упереджень. Однак емпіричні дослідження демонструють протилежне: інтелект не захищає від самообману, а часто робить його більш витонченим і стійким.

Цей парадокс пояснюється подвійною функцією мислення: воно служить не лише пізнанню, а також підтриманню когнітивної та ідентифікаційної узгодженості — підтриманню цілісності «Я» та зниженню внутрішньої напруги.

\section{Теоретичні передумови}

\subsection{Емпіричні дослідження}
Експериментальні та кореляційні дослідження мотивованого мислення (motivated reasoning) показують, що люди формують аргументацію під бажаний висновок. Когнітивні ресурси спрямовуються на пошук і конструювання підтверджуючих доказів.

Роботи \textbf{Dan Kahan \cite{kahan2013}} у рамках теорії культурної когніції (cultural cognition) виявили парадоксальний ефект: індивіди з вищим рівнем наукової грамотності та когнітивної рефлексії демонструють \textbf{сильнішу поляризацію} щодо суспільно значущих питань (кліматичні зміни, контроль зброї, гендерна політика). дані опитувань та лабораторних досліджень підтверджують, що більш когнітивно розвинені люди ефективніше захищають «свої» переконання за допомогою складнішої аргументації.

Класична теорія \textbf{когнітивного дисонансу Leon Festinger}, ефект \textbf{«ефект «сліпоти до власних упереджень» (bias blind spot)»} (Emily Pronin \cite{pronin2007}) та роботи \textbf{Daniel Kahneman \cite{kahneman2011}} (де аналітичне мислення Системи 2 часто виконує роль адвоката інтуїтивних висновків) у сукупності підтримують ключову тезу: інтелект не усуває упередження, а підвищує здатність до їх ефективної раціоналізації.

\textbf{Теоретичне твердження}: чим вищий когнітивний ресурс індивіда, тим більша його здатність до стабілізації переконань за умов мотивованого мислення.

\subsection{Когнітивний дисонанс}
Ці емпіричні результати добре узгоджуються з класичними моделями когнітивного дисонансу. Когнітивний дисонанс виникає при суперечності між переконаннями, діями та самооцінкою. Замість перегляду позиції індивід часто змінює інтерпретацію подій, їх значущість або окремі переконання, щоб відновити внутрішню гармонію.

\subsection{Мотивоване мислення}
У мотивованому мисленні висновок формується першим, а аргументація підбирається заднім числом. Когнітивні ресурси витрачаються не на пошук істини, а на підтвердження бажаного результату.

\subsection{Захисні механізми}
Серед психоаналітичних механізмів захисту ключовими є:
\begin{itemize}
    \item \textbf{Раціоналізація} — постфактум логічне виправдання поведінки чи позиції.
    \item \textbf{Інтелектуалізація} — відхід у абстрактне, теоретичне мислення для уникнення емоційного конфлікту.
\end{itemize}

\section{Динаміка трансформації переконань}
Зміна позиції часто відбувається не через поступовий аналіз нових фактів, а як реакція на внутрішній конфлікт. Типова послідовність:
\begin{enumerate}[label=\arabic*.]
    \item Виникнення екзистенційного або морального дискомфорту.
    \item Неможливість інтегрувати новий досвід у стару систему переконань.
    \item Формування нового висновку, який знижує напругу.
    \item Побудова складної аргументації під цей висновок.
\end{enumerate}

Таким чином, аргументація є вторинною щодо емоційно комфортної позиції.

\section{Роль інтелектуалізації}
Інтелектуалізація дозволяє дистанціюватися від неприємних емоцій і надати бажаній позиції вигляд об’єктивності та глибини. Використання складних філософських конструкцій підвищує суб’єктивну переконливість, ускладнює зовнішню критику та створює ілюзію рефлексивності. Особливо зручними стають концепції, які самі передбачають множинність інтерпретацій.

\section{Інверсія теоретичних концепцій}
У таких випадках теоретична складність перестає бути інструментом аналізу і перетворюється на механізм уникнення перевірки.
Концепції, створені як інструменти критичного аналізу (наприклад, деконструкція, множинність наративів, критика влади знання), можуть інвертуватися і перетворюватися на щит від критики. Ідея «будь-яка позиція є допустимою» в такому випадку перестає бути аналітичним інструментом і стає механізмом нейтралізації незручних аргументів.

\section{Селективна рефлексія та метакогнітивний розрив}
Індивід зберігає високу здатність до складного аналізу, але застосовує її селективно: гостро критикує зовнішні об’єкти та опонентів, водночас ігноруючи власні передумови. Це створює \textbf{метакогнітивний розрив} — рефлексія перестає рефлексувати саму себе.

\section{Поведінкові індикатори}
До типових поведінкових індикаторів належать: зміна стандартів доказовості залежно від того, чия позиція розглядається, ігнорування або применшення контрприкладів, уникнення прямої відповіді через зміну рамки дискусії та апеляцію до «складності питання» як до аргументу.

\section{Історичні та сучасні приклади}

\subsection{XX століття}
\textbf{Martin Heidegger} інтерпретував націонал-соціалістичний режим у категоріях історичної долі та онтологічного покликання. \textbf{Jean-Paul Sartre} у 1940--1950-х роках займав позиції, які частково легітимували радянський проєкт, інтерпретуючи його репресивні аспекти як історично зумовлені; після радянського вторгнення в Угорщину 1956 року він різко засудив дії СРСР і розірвав стосунки з французькою комуністичною партією.

У радянському інтелектуальному середовищі поширеною була раціоналізація через аргументи «історичної необхідності», «відсутності альтернатив» і «тимчасових жертв заради майбутнього».

\subsection{Сучасний контекст}
Дослідження Dan Kahan \cite{kahan2013} показують, що високоосвічені індивіди демонструють найсильнішу поляризацію щодо клімату, гендерних питань чи свободи слова. Складні теоретичні конструкції (критична теорія, поняття «безпеки» чи «влади дискурсу») часто використовуються для легітимації виключення опонентів або захисту власної позиції в академічних дебатах.

\section{Межі запропонованої моделі}
Запропонована модель не претендує на універсальне пояснення всіх форм зміни переконань. Вона фокусується на випадках, коли мислення виконує переважно захисну регулятивну функцію. У ситуаціях сильної мотивації до точності (наприклад, у професійній науковій діяльності за наявності жорстких інституційних перевірок) ефект інтелектуалізації та раціоналізації може бути суттєво ослаблений.

\section{Обмеження раціональної критики}
Зовнішня логічна критика є структурно обмеженою, оскільки стикається не з помилкою в аргументації, а із захисною функцією цих аргументів. Зміна позиції можлива переважно тоді, коли внутрішня «вартість» перегляду самообразу знижується.

\section{Висновки}
У цьому сенсі інтелект можна розглядати не як інструмент досягнення істини, а як механізм оптимізації когнітивної узгодженості.

Таким чином, інтелект виступає не лише як інструмент пізнання, а як механізм когнітивної стабілізації ідентичності.

Інтелект не є гарантією епістемічної чесності. Навпаки, він може бути потужним інструментом раціоналізації, стабілізації ідентичності та уникнення внутрішнього конфлікту.

Найвитонченіший самообман, як правило, має найкращу аргументацію. Ключовою умовою збереження пізнавальної функції мислення залишається радикальна здатність застосовувати рефлексію до власних передумов.

\textbf{Ключове формулювання}  
Інтелект не виправляє самообман — він його оптимізує.  
Рефлексія може перестати рефлексувати саму себе.

\bibliographystyle{plain}
\bibliography{refs/bibliography}

\end{document}
